\documentclass[12pt]{article}
% Préambule commun à tous les documents extraits de « La Revue CAPES »
\usepackage[T1]{fontenc}
\usepackage[utf8]{inputenc}
\usepackage{lmodern}
\usepackage{amsmath,amssymb,amsthm,amsfonts,mathrsfs}
\usepackage{setspace}
\usepackage{listings}
\usepackage{pgfplots}
\pgfplotsset{compat=1.18}
\usepackage{longtable}
\usepackage{adjustbox}
\usepackage{lettrine}
\usepackage{tkz-tab}
\usepackage{changepage}
\usepackage{pdflscape}
\usepackage{graphicx}
\usepackage{tikz}
\usepackage{xcolor}
\usepackage[a4paper,top=2cm,bottom=2.5cm,left=2.5cm,right=2cm]{geometry}
\usepackage{fancyhdr}
\usepackage{draftwatermark}
\SetWatermarkText{La RevueCapes}
\SetWatermarkLightness{0.9}
\SetWatermarkScale{2}
\usepackage{hyperref}
\hypersetup{colorlinks=true,linkcolor=blue!50!black,urlcolor=blue!50!black}

\definecolor{revue}{RGB}{20,60,120}

\pagestyle{fancy}
\fancyhf{}
\renewcommand{\headrulewidth}{0.4pt}
\setlength{\headheight}{15pt}

% \entete{Matière}{Titre}{Type de document}
\newcommand{\entete}[3]{%
  \fancyhead[L]{\small\textsc{La Revue CAPES} -- #1}%
  \fancyhead[R]{\small #2}%
  \fancyfoot[C]{\thepage}%
  \thispagestyle{plain}%
  \begin{center}
    {\color{revue}\rule{\textwidth}{1.2pt}}\\[4pt]
    {\small\textsc{Concours directs CAP-CEG / CAPES -- Burkina Faso}}\\[6pt]
    {\LARGE\bfseries\color{revue} #2}\\[6pt]
    {\large\itshape #1 -- #3}\\[2pt]
    {\color{revue}\rule{\textwidth}{1.2pt}}
  \end{center}
  \vspace{0.5cm}%
}

% Encadré pour les remarques ajoutées dans les corrigés
\newcommand{\remarque}[1]{\par\medskip\noindent\fcolorbox{revue}{revue!6}{\parbox{\dimexpr\linewidth-2\fboxsep-2\fboxrule}{\textbf{\color{revue}Remarque.} #1}}\par\medskip}


\begin{document}
\entete{Culture générale}{Corrigé du sujet type de dissertation n°15}{Correction}
\begin{spacing}{1.5}

\quad

\textbf{Introduction}

La Communauté économique des États de l’Afrique de l’Ouest (CEDEAO) s’est imposée comme un acteur clé dans la gestion des crises politiques et sécuritaires en Afrique de l’Ouest. Cependant, la récurrence des coups d’État, notamment au Burkina Faso, au Mali et en Guinée, soulève des interrogations sur l’efficacité de ses interventions. Entre l’imposition de sanctions économiques et politiques et son incapacité à prévenir la prolifération des putschs, la CEDEAO semble dépassée par les défis structurels de la région. Dès lors, comment expliquer son impuissance et quelles alternatives peuvent être envisagées pour une gestion plus efficace des crises en Afrique de l’Ouest ?

\quad

\textbf{Les sanctions de la CEDEAO : efficacité limitée}

La CEDEAO a souvent recours à des sanctions politiques et économiques pour répondre aux coups d’État. Par exemple, suite au putsch au Burkina Faso, des mesures restrictives ont été imposées, allant de la suspension du pays des organes de décision à des blocages financiers. L’objectif de ces sanctions est de rétablir l’ordre constitutionnel en exerçant une pression sur les régimes militaires.

Cependant, l’efficacité de ces mesures est souvent remise en question. Loin de ramener la stabilité, elles aggravent parfois les conditions de vie des populations locales, alimentant ainsi le ressentiment envers l’organisation régionale. Les sanctions renforcent également l’isolement des pays concernés, ce qui peut pousser les gouvernements militaires à chercher des alliances alternatives, notamment avec des puissances non occidentales comme la Russie.

\quad

\textbf{Une impuissance structurelle face aux coups d'État}

L’incapacité de la CEDEAO à prévenir les coups d’État découle de plusieurs facteurs structurels. Tout d’abord, les causes profondes des putschs, telles que la mauvaise gouvernance, la corruption, les inégalités sociales et la faiblesse des institutions démocratiques, ne sont pas directement abordées par l’organisation. Cela limite l’impact de ses interventions, qui se concentrent souvent sur les conséquences immédiates plutôt que sur les racines du problème.

De plus, la CEDEAO souffre d’un déficit de légitimité auprès des populations ouest-africaines, qui perçoivent parfois ses actions comme étant influencées par des intérêts étrangers ou élitistes. Cette perception affaiblit son rôle de médiateur et réduit son influence sur les gouvernements de transition. Enfin, les divisions internes entre les États membres, chacun ayant ses propres priorités et alliances, compliquent la prise de décisions cohérentes et efficaces.

\quad

\textbf{Perspectives pour une gestion plus efficace des crises}

Pour restaurer son efficacité, la CEDEAO doit adopter une approche plus proactive et inclusive. Cela implique de s’attaquer aux causes profondes des instabilités politiques, notamment à travers le soutien à la bonne gouvernance, le renforcement des institutions démocratiques et la promotion d’un développement économique équitable.

En outre, il est essentiel que l’organisation renforce sa légitimité en impliquant davantage la société civile et les populations locales dans ses processus décisionnels. Une approche plus participative pourrait réduire les tensions et encourager des solutions adaptées aux réalités locales. Par ailleurs, la CEDEAO devrait également collaborer étroitement avec d’autres organisations continentales, telles que l’Union africaine, pour coordonner ses efforts et mutualiser les ressources.

\quad

\textbf{Conclusion}

Face à la crise burkinabè et à la prolifération des coups d’État en Afrique de l’Ouest, la CEDEAO se retrouve à un carrefour. Si ses sanctions politiques et économiques ont montré leurs limites, elles soulignent également la nécessité d’une refonte stratégique de son mode d’intervention. Pour rester un acteur pertinent, la CEDEAO devra renforcer sa légitimité, adopter une approche holistique et inclusive, et s’attaquer aux causes profondes des instabilités politiques dans la région. L’avenir de la stabilité en Afrique de l’Ouest dépendra largement de sa capacité à s’adapter aux défis complexes de la région.



\quad

\end{spacing}
\end{document}
